% ============================================================
% Section: Experiments
% ============================================================

\section{Experiments}
\label{sec:experiments}

% ============================================================
\subsection{Experimental Setup}
\label{sec:exp:setup}

% TODO: Describe experimental setup.
%   - Dataset splits (ITB Benchmark + ISIC \cite{codella2018skin,codella2019skin})
%   - Implementation details: GPU, batch size, optimizer, LR schedule, epochs
%   - Evaluation metrics: PLCC, SRCC, ECE, NLL, PSNR, SSIM, LPIPS, AUC

\todo{Write experimental setup here.}

% ============================================================
\subsection{Quality Assessment Results}
\label{sec:exp:iqa}

% TODO: Compare VisiScore-Net against NR-IQA baselines.
%   Table: PLCC / SRCC on ITB test set vs. BRISQUE, NIQE, MUSIQ, HyperIQA
%   Analysis: domain-specific gain, low-quality tail performance.

\todo{Write quality assessment results here.}

% \input{tables/tab_iqa.tex}

% ============================================================
\subsection{Uncertainty Estimation Results}
\label{sec:exp:uncertainty}

% TODO: Compare Q-VIB uncertainty quality.
%   Table: ECE / NLL vs. MC Dropout \cite{gal2016dropout},
%          Deep Ensembles \cite{lakshminarayanan2017simple},
%          Temperature Scaling \cite{guo2017calibration}
%   Figure: reliability diagram, uncertainty-quality correlation scatter plot.

\todo{Write uncertainty estimation results here.}

% \input{tables/tab_uncertainty.tex}

% ============================================================
\subsection{Enhancement Results}
\label{sec:exp:enhance}

% TODO: Compare VisiEnhance-Net against restoration baselines.
%   Table: PSNR / SSIM / LPIPS vs. NAFNet \cite{chen2022simple},
%          Restormer \cite{zamir2022restormer}, blind enhancement baselines.
%   Qualitative: figure grid (input / NAFNet / Restormer / ours / GT).

\todo{Write enhancement results here.}

% \input{tables/tab_enhancement.tex}

% ============================================================
\subsection{End-to-End Diagnostic Evaluation}
\label{sec:exp:downstream}

% TODO: Show that the dual-channel pipeline improves downstream diagnosis.
%   Table: AUC / sensitivity / specificity on ISIC test set,
%          with/without enhancement, with/without uncertainty routing.
%   Analysis: benefit is concentrated in low-quality subset.

\todo{Write end-to-end diagnostic evaluation here.}

% ============================================================
\subsection{Ablation Study}
\label{sec:exp:ablation}

% TODO: Ablation table with following variants:
%   (a) VisiScore-Net only (no uncertainty)
%   (b) Q-VIB fixed beta (no quality conditioning)
%   (c) VisiEnhance-Net without FiLM conditioning
%   (d) Agent without uncertainty routing
%   (e) Full model (ours)

\todo{Write ablation study here.}

% \input{tables/tab_ablation.tex}

% ============================================================
\subsection{Analysis and Discussion}
\label{sec:exp:discussion}

% TODO:
%   - Failure case analysis (very dark lesions, hair occlusion)
%   - Computational cost comparison (FLOPs, inference time)
%   - Sensitivity analysis on quality threshold tau

\todo{Write analysis and discussion here.}
